\subsection{Общие сведения}

Рассматриваемая система аналого-цифрового преобразования и цифровой обработки сигналов с воздействием на объект управления посредством широтно-импульсной модуляции (рис.~\eqref{fig: discrete_time_system}).

\begin{figure}
	\center
	\includegraphics [scale=1.00] {frequencyResponse/discreteTimeSystem.drawio.png}

	\caption{Рассматриваемая система}
	\label{fig: discrete_time_system}
\end{figure}

Операции выполняются циклически с постоянным периодом ${T_{d}}$, соответствующим базовой частоте дискретного регулирования ${f_{d} = 1 / T_{d}}$.

В канал управления системы аддитивно инжектируется сигнал в виде дискретной последовательности ${z[n]}$.

Откликом системы являются два континуальных сигнала: ${x_{c}(t)}$, ${y_{c}(t)}$.

Система преобразует континуальные сигналы ${x_{c}(t)}$, ${y_{c}(t)}$ в дискретные последовательности ${x[n]}$, ${y[n]}$.

Преобразование осуществляется посредством дискретизации:
\begin{equation}\label{eqn: sampleRule}
    \begin{split}
        &x[n] = x_{c} (n T_s), \\
        &y[n] = y_{c} (n T_s),
    \end{split}
\end{equation}
где ${T_{s}}$ -- период дискретизации, ${T_{s} \geq T_{d}}$.

Коэффициенты разложения дискретных последовательностей
\begin{equation*}
    \begin{aligned}
        &x\left[n\right] = \cfrac{1}{N} \sum_{l=0}^{N - 1} {X\left[l\right] e^{j\cfrac{2\pi}{N}ln}}, \\
        &y\left[n\right] = \cfrac{1}{N} \sum_{l=0}^{N - 1} {Y\left[l\right] e^{j\cfrac{2\pi}{N}ln}}
    \end{aligned}
\end{equation*}
определяются по $N$-последовательным отсчетам дискретных последовательностей ${x[n]}$, ${y[n]}$:
\begin{equation}\label{eqn: DFT}
    \begin{aligned}
        &X\left[l\right] = \sum_{n=0}^{N - 1} {x\left[n\right] e^{-j\cfrac{2\pi}{N}nl}}, \\
        &Y\left[l\right] = \sum_{n=0}^{N - 1} {y\left[n\right] e^{-j\cfrac{2\pi}{N}nl}}.
    \end{aligned}
\end{equation}

Коэффициенты ${X\left[l\right], Y\left[l\right] \in \mathbb{C}}$, представляют собой дискретные периодические последовательности с периодом $N$.

При условии, что измерительное окно длительностью ${T_{h}}$ вмещает целое количество $N$ периодов дискретизации ${T_{s}}$ (${T_{h} = N \cdot T_{s}}$) дискретные последовательности ${X\left[l\right]}$, ${Y\left[l\right]}$ на отрезке ${l \in \left[0; N - 1 \right]}$ длины $N$ содержат коэффициенты разложения ${x\left[n\right]}$, ${y\left[n\right]}$ в диапазоне частот ${f \in \left[0, f_{s} / 2\right]}$ с шагом по частоте ${\Delta f=1 / T_{h}}$.

Составляющим с частотой ${f_{h}}$ соответствуют ${X\left[l\right]}$, ${Y\left[l\right]}$ при ${l = 1}$, а с частотой ${f_{s} / 2}$ -- при ${l = \bigl \lfloor (N - 1)/2 \bigr \rfloor}$.

В данном разделе рассматриваются алгоритмы измерения частотных характеристик, основанные на воздействии на систему дискретной последовательностью ${z[n]}$ моногармонического типа, определяемой косинусоидальной функцией
\begin{equation*}
    z[n] = Z \cos{(\underbrace{2\pi f_{z}}_{\omega_{z}} T_{d} n )}
\end{equation*}
с частотой ${f_{z}}$.

Минимальное количество отсчетов $N$ на измерительном окне ${T_{h}}$ принято равным 3, поэтому частота ${f_{z}}$ ограничена сверху значением ${f_{d} / 3}$.

Модель сигналов-откликов системы принята следующей:
\begin{equation*}
    \begin{aligned}
        &x_c(t) &= X_0 &+ X_{z} \cos{\left(\omega_{z} t + \phi_{x} \right)} &+ X_{g} \cos{\left(\omega_{g} t\right)}, \\
        &y_c(t) &= Y_0 &+ Y_{z} \cos{\left(\omega_{z} t + \phi_{y} \right)} &+ Y_{g} \cos{\left(\omega_{g} t\right)},
    \end{aligned}
\end{equation*}
где ${\omega_{z} = 2\pi f_{z}}$.

Алгоритм измерения частотной характеристики звена системы основан на последовательной смене периода инжектируемого сигнала ${T_{z} = 1 / f_{z}}$ в серии из $K$-измерений.

Каждому $k$-му измерению, ${k \in \left[1, K\right]}$, соответствует пара параметров $Z^{(k)}$, ${T^{(k)}_{z}}$, следовательно, и набор дискретных последовательностей ${x^{(k)}[n]}$, ${y^{(k)}[n]}$, ${z^{(k)}[n]}$, а также коэффициентов ${X^{(k)}\left[l\right]}$, ${Y^{(k)}\left[l\right]}$.

Под частотной характеристикой понимается набор пар значений ${f^{(k)}_{z}}$, ${W^{(k)}}$, где в каждом $k$-ом измерении
\begin{equation}\label{eqn: complexFreqResp}
    W^{(k)} = \cfrac{Y^{(k)}\left[L^{(k)}\right]}{X^{(k)}\left[L^{(k)}\right]} = \left| W^{(k)} \right| e^{j \angle W^{(k)}},
\end{equation}
${L^{(k)} = {T^{(k)}_{h}} / {T^{(k)}_{z}}}$.

Пары значений ${f^{(k)}_{z}}$, ${\left| W^{(k)} \right|}$ и ${f^{(k)}_{z}}$, ${\angle W^{(k)}}$ образуют амплитудно-частотную и фазо-частотную характеристики, соответственно.
\begin{equation}\label{eqn: abs_angleFreqResp}
    \begin{aligned}
        &\left| W^{(k)} \right| = \cfrac{\left| Y^{(k)}\left[L^{(k)}\right] \right| }{ \left| X^{(k)}\left[L^{(k)}\right] \right|}, \\
        &\angle W^{(k)} = \angle Y^{(k)}\left[L^{(k)}\right] - \angle X^{(k)}\left[L^{(k)}\right].
    \end{aligned}
\end{equation}